% Generated draft: requires mathematical and source review.
\documentclass[11pt]{article}
\usepackage[margin=25mm]{geometry}
\usepackage{fontspec}
\setmainfont{Latin Modern Roman}
\usepackage{amsmath,amssymb,mathtools,unicode-math}
\setmathfont{Latin Modern Math}
\usepackage{enumitem,needspace,xurl,hyperref}
\setlength{\parindent}{0pt}
\setlength{\parskip}{5pt}
\setlength{\emergencystretch}{3em}
\newcommand{\N}{\mathbb{N}}
\newcommand{\Z}{\mathbb{Z}}
\newcommand{\Q}{\mathbb{Q}}
\newcommand{\R}{\mathbb{R}}
\newcommand{\C}{\mathbb{C}}
\newcommand{\F}{\mathbb{F}}
\newcommand{\A}{\mathbb{A}}
\newcommand{\PP}{\mathbb{P}}
\newcommand{\E}{\mathbb{E}}
\newcommand{\eps}{\varepsilon}
\newcommand{\dd}{\,\mathrm{d}}
\begin{document}
\section*{UnsolvedMath ID: 11000054}
\textbf{Source problem code:} catalogue entry
\subsection{Definitions and mathematical statement}
Let \(\mathcal I_g\) be the Torelli group in the source's surface convention, the kernel of the mapping-class action on integral first homology. The excerpt does not specify closed versus boundary or marked-point conventions; these matter for sharp generator counts.

\emph{Problem 5.7 (cubic generating-set problem).} Find a family of generating sets \(S_g\) for \(\mathcal I_g\) of cardinality \(O(g^d)\) for some fixed \(d\geq3\), optimally \(d=3\). Explicitly, seek a construction and constants \(C,g_0\), independent of \(g\), with
\[
 \langle S_g\rangle=\mathcal I_g,\qquad |S_g|\leq Cg^d
 \quad(g\geq g_0)
\]
in the intended genus range. The polynomial exponent must not be allowed to depend on \(g\).

The source separately quotes Johnson's harder question from [Jo2], \(\S5\): for every \(g\geq4\), can \(\mathcal I_g\) be generated by
\[
 \frac{4g^3-g}{3}
\]
elements? It says this would be sharp, with the lower-bound context omitted. It reports Johnson [Jo2], Theorem 3, constructing \(35\) explicit generators in genus \(3\), attaining the corresponding sharp count. It describes converting his earlier generators to cubic size as unwieldy for \(g>3\). The extracted earlier bound reads \(O(2g)\); in this context an exponential \(O(2^g)\) is a plausible lost-superscript reading, but the typography is unresolved. It is not silently read as a linear-size set, which would already satisfy the polynomial request.

The proposed approach follows Johnson's plan using a simpler generating set for \(\operatorname{Mod}_g\). The source credits Brendle and the author [BFa1] with six involutions, and Kassabov [Ka] with four involutions for at least \(g\geq7\). It explains that two involutions would give a quotient of the infinite dihedral group and hence a virtually abelian group, precluding that possibility for the intended mapping class groups. It then asks naturally for sharp involution-generator bounds. These are related historical facts and motivations, with source scope, not a claimed present-day optimum or a solution to the Torelli problem.
\subsection{Short English statement}
Construct polynomial-size generating sets for Torelli groups, ideally cubic in genus. Johnson's separate sharp target is \((4g^3-g)/3\) for \(g\geq4\), with \(35\) generators reported in genus three; surface conventions and one flattened exponential bound require checking.
\subsection{Sources}
\begin{itemize}
\item Local UnsolvedMath registry, permanent ID 11000054, source code catalogue entry, Problem 5.7, ``Cubic genset problem,'' PDF page 45. No individual proposer is supplied for the broad problem; Johnson is credited for the sharper question.
\item Farb, \emph{Problems on Mapping Class Groups and Related Topics}. Supplied URL: \url{https://www.math.uchicago.edu/~farb/papers/mcgbook.pdf}; recorded access 2026-07-29; extraction pdf-font-structure. The edited volume contains 329 explicitly typeset headings across chapters, with repeated local numbering disambiguated by page.
\item Johnson [Jo2], \(\S5\) and Theorem 3; Brendle and the author [BFa1]; Kassabov [Ka]. Full bibliographic entries are not supplied. The source states the six- and four-involution results and their quoted genus restriction.
\item Supplied status open and OPEN-TRIAGE. The summary says only ``Cannot verify.'' Historical current-bound language in the source is not treated as a verified 2026 literature assessment. Current primary-literature status and redistribution terms remain NEEDS\_REVIEW.
\end{itemize}
\end{document}
